% from ega4-17-fr.tex
\oldpage[IV-4]{80}

où $Y=\Spec(A)$, $A$ étant un anneau local. Comme $f$ est de présentation finie, on a $X=\Spec(B)$, où $B$ est un $A$-module plat de présentation finie (1.4.7), donc \emph{libre} (Bourbaki, \emph{Alg. comm.}, chap. II, § 5, n\textsuperscript{o} 2, cor. 2 du th. 1). En outre, si $\mathfrak m$ est l'idéal maximal de $A$ et $k$ son corps résiduel, $B/\mathfrak mB$ est par hypothèse un corps, à la fois extension radicielle et extension finie séparable de $k$, puisque $f$ est étale et radiciel (17.6.1); donc $B/\mathfrak mB$ est isomorphe à $k$, et comme $B$ est un $A$-module libre, $B$ est isomorphe à $A$. C.Q.F.D.

\begin{corollary}[17.9.2]
\label{IV.17.9.2}
Soit $X$ un préschéma connexe; si $f:X\to Y$ est une immersion fermée étale, alors $f$ est un isomorphisme de $X$ sur une composante connexe ouverte de $Y$.
\end{corollary}

En effet $f$ est étale et radiciel, donc un isomorphisme de $X$ sur un préschéma induit sur une partie ouverte de $Y$; mais par hypothèse $f(X)$ est fermé dans $Y$, donc à la fois ouvert et fermé, et puisque $f(X)$ est connexe, c'est une composante connexe de $Y$.

\begin{corollary}[17.9.3]
\label{IV.17.9.3}
Soit $f:X\to Y$ un morphisme étale (resp. étale et séparé). Alors toute $Y$-section $g:Y\to X$ de $X$ est une immersion ouverte (resp. ouverte et fermée); en outre l'application $g\mapsto g(Y)$ est une bijection de l'ensemble $\Gamma(X/Y)$ des $Y$-sections de $X$ sur l'ensemble des parties ouvertes $Z$ (resp. ouvertes et fermées) de $X$ telles que la restriction de $f$ à $Z$ soit un morphisme surjectif et radiciel de $Z$ sur $Y$.
\end{corollary}

En effet, le fait que $g$ soit une immersion ouverte résulte déjà de ce que $f$ est non ramifié (17.4.1, $b''$)), et la restriction de $f$ à l'ouvert $g(Y)$ de $X$ est un isomorphisme. Inversement, si $Z$ est un ouvert de $X$ tel que $f|Z$ soit un morphisme surjectif et radiciel de $Z$ sur $Y$, $f|Z$ est un isomorphisme en vertu de (17.9.1), puisqu'il est étale. Si $f$ est en outre séparé, on sait que $g$ est une immersion fermée (1, 5.4.6), ce qui achève de prouver le corollaire.

\begin{corollary}[17.9.4]
\label{IV.17.9.4}
Soient $Y$ un préschéma connexe, $f:X\to Y$ un morphisme étale et séparé. Alors toute $Y$-section $g$ de $X$ est un isomorphisme de $Y$ sur une composante connexe ouverte de $X$, et l'application $g\mapsto g(Y)$ est une bijection de $\Gamma(X/Y)$ sur l'ensemble des composantes connexes ouvertes $Z$ de $X$ telles que la restriction de $f$ à $Z$ soit un morphisme surjectif et radiciel de $Z$ sur $Y$.
\end{corollary}

\begin{corollary}[17.9.5]
\label{IV.17.9.5}
Soient $g:Y\to S$, $h:X\to S$ deux morphismes localement de présentation finie; on suppose en outre que $h$ est plat. Pour qu'un $S$-morphisme $f:X\to Y$ soit une immersion ouverte (resp. un isomorphisme), il faut et il suffit que, pour tout $s\in S$, le morphisme $f_s:h^{-1}(s)\to g^{-1}(s)$ déduit de $f$ par le changement de base $\Spec(k(s))\to S$, soit une immersion ouverte (resp. un isomorphisme).
\end{corollary}

En effet, si $f_s$ est une immersion ouverte pour tout $s\in S$, il résulte de (17.8.3) que $f$ est un morphisme étale; comme pour tout $y\in Y$, on a $f^{-1}(y)=f_s^{-1}(y)$ avec $s=g(y)$, $f$ est radiciel; donc $f$ est une immersion ouverte en vertu de (17.9.1). Si de plus $f_s$ est surjectif pour tout $s\in S$, $f$ est surjectif, donc un isomorphisme.

La proposition suivante précise (10.4.11):

\begin{proposition}[17.9.6]
\EGAReaderTarget{ega4.c04.prop.17-9-6}{IV.17.9.6}
Soient $S$ un préschéma, $X$ un $S$-préschéma de présentation finie. Tout $S$-endomorphisme de $X$ qui est un monomorphisme est un automorphisme de $X$.
\end{proposition}

Soient $f:X\to S$ le morphisme structural, $g$ le $S$-endomorphisme considéré. La question est locale sur $S$, et on peut par suite supposer que $S$ est affine. Utilisant (8.9.1) et (8.10.5, (i \emph{bis})), on est ramené au cas où $S=\Spec(A)$, où $A$ est une $\mathbb{Z}$-algèbre de type fini, et par suite $X$ est un $\mathbb{Z}$-préschéma de type fini. Il résulte déjà de (10.4.11) que $g$ est un morphisme bijectif, puisqu'un monomorphisme est radiciel (8.11.5.1); il suffira donc de montrer

